\begin{references}
\new{Achard, J. A., Gajdos Preuss, T., \& Servant, M. (2025).} Extending continuous flow models of immediate decision reports to delayed decision reports. \emph{Journal of Experimental Psychology. General}. https://doi.org/10.1037/xge0001748

\del{Afacan-Seref, K., Steinemann, N. A., Blangero, A., \& Kelly, S. P. (2018). Dynamic Interplay of Value and Sensory Information in High-Speed Decision Making.} \del{\emph{Current Biology: CB}}\del{,} \del{\emph{28}}\del{(5), 795-802.e6. https://doi.org/10.1016/j.cub.2018.01.071}

\del{Aza\"is, J.-M., \& Wschebor, M. (2009). Level sets and extrema of random processes and fields. John Wiley \& Sons. https://doi.org/10.1002/9780470434642}

Balsdon, T., Verdonck, S., Loossens, T., \& Philiastides, M. G. (2023). Secondary motor integration as a final arbiter in sensorimotor decision-making. \emph{PLoS Biology}, \emph{21}(7), e3002200. https://doi.org/10.1371/journal.pbio.3002200

Burle, B., Tandonnet, C., \& Hasbroucq, T. (2010). Excitatory and inhibitory motor mechanisms of temporal preparation. In A. C. Nobre \& J. T. Coull (Eds.), \emph{Attention and Time} (pp. 243--256). Oxford University Press. https://doi.org/10.1093/acprof:oso/9780199563456.003.0018

Churchland, M. M., \& Shenoy, K. V. (2024). Preparatory activity and the expansive null-space. \emph{Nature Reviews. Neuroscience}, \emph{25}(4), 213--236. https://doi.org/10.1038/s41583-024-00796-z

\new{Cisek, P., \& Kalaska, J. F. (2005). Neural correlates of reaching decisions in dorsal premotor cortex: Specification of multiple direction choices and final selection of action.} \new{\emph{Neuron}}\new{,} \new{\emph{45}}\new{(5), 801--814. https://doi.org/10.1016/j.neuron.2005.01.027}

Coles, M. G., Gratton, G., Bashore, T. R., Eriksen, C. W., \& Donchin, E. (1985). A psychophysiological investigation of the continuous flow model of human information processing. \emph{Journal of Experimental Psychology. Human Perception and Performance}, \emph{11}(5), 529--553. https://doi.org/10.1037//0096-1523.11.5.529

\del{Corbett, E. A., Martinez-Rodriguez, L. A., Judd, C., O'Connell, R. G., \& Kelly, S. P. (2023). Multiphasic value biases in fast-paced decisions.} \del{\emph{eLife}}\del{,} \del{\emph{12}}\del{, e67711. https://doi.org/10.7554/eLife.67711}

Dendauw, E., Evans, N. J., Logan, G. D., Haffen, E., Bennabi, D., Gajdos, T., \& Servant, M. (2024). The gated cascade diffusion model: An integrated theory of decision making, motor preparation, and motor execution. \emph{Psychological Review}. https://doi.org/10.1037/rev0000464

Dendauw, E., Logan, G. D., Schall, J. D., Gajdos Preuss, T., \& Servant, M. (2026). Separating decision and motor contributions to behavioral biases induced by manipulating stimulus probability. \emph{Cognitive Psychology}, \emph{165}, 101818. https://doi.org/10.1016/j.cogpsych.2026.101818

Donner, T. H., Siegel, M., Fries, P., \& Engel, A. K. (2009). Buildup of Choice-Predictive Activity in Human Motor Cortex during Perceptual Decision Making. \emph{Current Biology}, \emph{19}(18), 1581--1585. https://doi.org/10.1016/j.cub.2009.07.066

Eimer, M. (1998). The lateralized readiness potential as an on-line measure of central response activiation processes. \emph{Behavior Research Methods, Instruments \& Computers}, \emph{30}(1), 146--156. https://doi.org/10.3758/BF03209424

Forstmann, B. U., Ratcliff, R., \& Wagenmakers, E.-J. (2016). Sequential Sampling Models in Cognitive Neuroscience: Advantages, Applications, and Extensions. \emph{Annual Review of Psychology}, \emph{67}, 641--666. https://doi.org/10.1146/annurev-psych-122414-033645

Gold, J. I., \& Shadlen, M. N. (2007). The neural basis of decision making. \emph{Annual Review of Neuroscience}, \emph{30}, 535--574. https://doi.org/10.1146/annurev.neuro.29.051605.113038

Gratton, G., Coles, M. G., Sirevaag, E. J., Eriksen, C. W., \& Donchin, E. (1988). Pre- and poststimulus activation of response channels: A psychophysiological analysis. \emph{Journal} \emph{of Experimental Psychology. Human Perception and Performance}, \emph{14}(3), 331--344. https://doi.org/10.1037//0096-1523.14.3.331

Heideman, S. G., van Ede, F., \& Nobre, A. C. (2018). Temporal alignment of anticipatory motor cortical beta lateralisation in hidden visual-motor sequences. \emph{The European Journal of Neuroscience}, \emph{48}(8), 2684--2695. https://doi.org/10.1111/ejn.13700

Inagaki, H. K., Chen, S., Daie, K., Finkelstein, A., Fontolan, L., Romani, S., \& Svoboda, K. (2022). Neural Algorithms and Circuits for Motor Planning. \emph{Annual Review of Neuroscience}, \emph{45}, 249--271. https://doi.org/10.1146/annurev-neuro-092021-121730

Kelly, S. P., Corbett, E. A., \& O'Connell, R. G. (2021). Neurocomputational mechanisms of prior-informed perceptual decision-making in humans. \emph{Nature Human Behaviour}, \emph{5}(4), 467--481. https://doi.org/10.1038/s41562-020-00967-9

Kelly, S. P., \& O'Connell, R. G. (2013). Internal and External Influences on the Rate of Sensory Evidence Accumulation in the Human Brain. \emph{Journal of Neuroscience}, \emph{33}(50), 19434--19441. https://doi.org/10.1523/JNEUROSCI.3355-13.2013

Kiani, R., Hanks, T. D., \& Shadlen, M. N. (2008). Bounded integration in parietal cortex underlies decisions even when viewing duration is dictated by the environment. \emph{The Journal of Neuroscience: The Official Journal of the Society for Neuroscience}, \emph{28}(12), 3017--3029. https://doi.org/10.1523/JNEUROSCI.4761-07.2008

Li, J., Zhai, M., Liu, L., Qian, Y., \& Liao, Y. (2026). Beta Oscillations as Internal Timers: Independent Contributions of Temporal Predictability to Motor Preparation. \emph{Psychophysiology}, \emph{63}(2), e70249. https://doi.org/10.1111/psyp.70249

Liu, J., \& Liu, Q. (2016). Use of the integrated profile for voluntary muscle activity detection using EMG signals with spurious background spikes: A study with incomplete spinal cord injury. \emph{Biomedical Signal Processing and Control}, \emph{24}, 19--24. https://doi.org/10.1016/j.bspc.2015.09.004

Mazurek, M. E., Roitman, J. D., Ditterich, J., \& Shadlen, M. N. (2003). A Role for Neural Integrators in Perceptual Decision Making. \emph{Cerebral Cortex}, \emph{13}(11), 1257--1269. https://doi.org/10.1093/cercor/bhg097

Morey, R. D. (2008). Confidence intervals from normalized data: A correction to Cousineau (2005). \emph{Tutorial in Quantitative Methods for Psychology}, \emph{4}(2), 61--64.

\del{Murphy, P. R., Boonstra, E., \& Nieuwenhuis, S. (2016). Global gain modulation generates time-dependent urgency during perceptual choice in humans.} \del{\emph{Nature Communications}}\del{,} \del{\emph{7}}\del{, 13526. https://doi.org/10.1038/ncomms13526}

O'Connell, R. G., Dockree, P. M., \& Kelly, S. P. (2012). A supramodal accumulation-to-bound signal that determines perceptual decisions in humans. \emph{Nature Neuroscience}, \emph{15}(12), 1729--1735. https://doi.org/10.1038/nn.3248

O'Connell, R. G., \& Kelly, S. P. (2021). Neurophysiology of Human Perceptual Decision-Making. \emph{Annual Review of Neuroscience}, \emph{44}, 495--516. https://doi.org/10.1146/annurev-neuro-092019-100200

Pfurtscheller, G., \& Lopes da Silva, F. H. (1999). Event-related EEG/MEG synchronization and desynchronization: Basic principles. \emph{Clinical Neurophysiology: Official Journal of the International Federation of Clinical Neurophysiology}, \emph{110}(11), 1842--1857. https://doi.org/10.1016/s1388-2457(99)00141-8

Ramayya, A. G., Buch, V., Richardson, A., Lucas, T., \& Gold, J. I. (2025). Human response times are governed by dual anticipatory processes with distinct neural signatures. \emph{Communications Biology}, \emph{8}(1), 124. https://doi.org/10.1038/s42003-025-07516-y

\new{Ratcliff, R., \& Tuerlinckx, F. (2002). Estimating parameters of the diffusion model: Approaches to dealing with contaminant reaction times and parameter variability. Psychonomic Bulletin \& Review, 9(3), 438--481. https://doi.org/10.3758/BF03196302}

Ratcliff, R., \& Van Dongen, H. P. A. (2011). Diffusion model for one-choice reaction-time tasks and the cognitive effects of sleep deprivation. \emph{Proceedings of the National Academy of Sciences of the United States of America}, \emph{108}(27), 11285--11290. https://doi.org/10.1073/pnas.1100483108

Resulaj, A., Kiani, R., Wolpert, D. M., \& Shadlen, M. N. (2009). Changes of mind in decision-making. \emph{Nature}, \emph{461}(7261), 263--266. https://doi.org/10.1038/nature08275

\new{Roitman, J. D., \& Shadlen, M. N. (2002). Response of neurons in the lateral intraparietal area during a combined visual discrimination reaction time task.} \emph{The Journal of Neuroscience: The Official Journal of the Society for Neuroscience}, \emph{22}\new{(21), 9475--9489. https://doi.org/10.1523/JNEUROSCI.22-21-09475.2002}

Rogge, J., Jocham, G., \& Ullsperger, M. (2022). Motor cortical signals reflecting decision making and action preparation. \emph{NeuroImage}, \emph{263}, 119667. https://doi.org/10.1016/j.neuroimage.2022.119667

Santello, M., \& McDonagh, M. J. (1998). The control of timing and amplitude of EMG activity in landing movements in humans. \emph{Experimental Physiology}, \emph{83}(6), 857--874. https://doi.org/10.1113/expphysiol.1998.sp004165

Selen, L. P. J., Shadlen, M. N., \& Wolpert, D. M. (2012). Deliberation in the motor system: Reflex gains track evolving evidence leading to a decision. \emph{The Journal of Neuroscience: The Official Journal of the Society for Neuroscience}, \emph{32}(7), 2276--2286. https://doi.org/10.1523/JNEUROSCI.5273-11.2012

Servant, M., Logan, G. D., Gajdos, T., \& Evans, N. J. (2021). An integrated theory of deciding and acting. \emph{Journal of Experimental Psychology. General}, \emph{150}(12), 2435--2454. https://doi.org/10.1037/xge0001063

Servant, M., Tillman, G., Schall, J. D., Logan, G. D., \& Palmeri, T. J. (2019). Neurally constrained modeling of speed-accuracy tradeoff during visual search: Gated accumulation of modulated evidence. \emph{Journal of Neurophysiology}, \emph{121}(4), 1300--1314. https://doi.org/10.1152/jn.00507.2018

Smith, P. L. (1995). Psychophysically principled models of visual simple reaction time. \emph{Psychological Review}, \emph{102}(3), 567--593. https://doi.org/10.1037/0033-295X.102.3.567

Twomey, D. M., Kelly, S. P., \& O'Connell, R. G. (2016). Abstract and Effector-Selective Decision Signals Exhibit Qualitatively Distinct Dynamics before Delayed Perceptual Reports. \emph{The Journal of Neuroscience: The Official Journal of the Society for Neuroscience}, \emph{36}(28), 7346--7352. https://doi.org/10.1523/JNEUROSCI.4162-15.2016

Usher, M., \& McClelland, J. L. (2001). The time course of perceptual choice: The leaky, competing accumulator model. \emph{Psychological Review}, \emph{108}(3), 550--592. https://doi.org/10.1037/0033-295x.108.3.550

\new{van Ede, F., \& Nobre, A. C. (2023). Turning attention inside out: How working memory serves behavior.} \emph{Annual Review of Psychology}, \emph{74}\new{, 137--165. https://doi.org/10.1146/annurev-psych-021422-041757}

\del{van Vugt, M. K., Simen, P., Nystrom, L., Holmes, P., \& Cohen, J. D. (2014). Lateralized Readiness Potentials Reveal Properties of a Neural Mechanism for Implementing a Decision Threshold.} \del{\emph{PLoS ONE}}\del{,} \del{\emph{9}}\del{(3), e90943. https://doi.org/10.1371/journal.pone.0090943}

Verdonck, S., Loossens, T., \& Philiastides, M. G. (2020). The Leaky Integrating Threshold and its impact on evidence accumulation models of choice response time (RT). \emph{Psychological Review}. https://doi.org/10.1037/rev0000258

Wang, X.-J. (2008). Decision making in recurrent neuronal circuits. \emph{Neuron}, \emph{60}(2), 215--234. https://doi.org/10.1016/j.neuron.2008.09.034

\end{references}
